\section{Revision: explain, compare, apply}
\subsection{The essential chain of reasoning}
\begin{enumerate}
\item A robot acts physically through actuators and obtains information through sensors.
\item Engineered, repetitive tasks can often use a prepared plan and a sufficiently complete model.
\item Unknowns and changes require current sensing, execution monitoring and adaptation.
\item Hierarchical organization supports explicit plans; reactive organization supports immediate responses.
\item Hybrid organization links these capabilities, while human--robot interaction adds control allocation and shared understanding.
\end{enumerate}

\subsection{Questions with concise model answers}
\begin{samepage}\textbf{1. Is every robot an intelligent robot?}\\
No. A robot may mainly replay a prepared sequence. An intelligent robot uses information about its situation to select or adapt actions toward a task goal.\par\end{samepage}

\begin{samepage}\textbf{2. Why does field robotics make perception important?}\\
Its environment cannot usually be arranged and maintained like a production fixture. The robot must observe relevant changes instead of relying only on a prepared description.\par\end{samepage}

\begin{samepage}\textbf{3. What separates telepresence from autonomy?}\\
Telepresence improves the human's experience of the remote environment. Autonomy concerns which decisions and tasks the robot can perform independently.\par\end{samepage}

\begin{samepage}\textbf{4. How do shared control and control trading differ in Murphy?}\\
Shared control involves assistance or monitored delegation with human involvement. Control trading delegates a task the robot can complete autonomously until a new command or interruption.\par\end{samepage}

\begin{samepage}\textbf{5. Explain automation versus autonomy without using robot shape.}\\
Automation emphasizes repeated execution under well-modeled conditions. Autonomy emphasizes adaptation when relevant conditions or outcomes are uncertain. Both can exist in the same robot.\par\end{samepage}

\begin{samepage}\textbf{6. What happens to an unlisted obstacle under CWA and OWA?}\\
Under CWA, the model treats it as absent. Under OWA, its absence from the model leaves its existence unknown. Current observations are needed to determine whether a route is clear.\par\end{samepage}

\begin{samepage}\textbf{7. Can closed-loop control be used under CWA?}\\
Yes. Feedback may correct a joint's position while the workspace and task are assumed complete and known. Model completeness and feedback are different properties.\par\end{samepage}

\begin{samepage}\textbf{8. Why can a valid symbolic plan fail?}\\
It is valid relative to its model. An omitted condition, an unexpected event or an action whose physical effect differs from its modeled effect can prevent execution.\par\end{samepage}

\begin{samepage}\textbf{9. What is the frame problem?}\\
Representing the effects of actions and the persistence of unaffected facts without an impractical enumeration. Murphy also emphasizes the broader burden of maintaining a tractable world model.\par\end{samepage}

\begin{samepage}\textbf{10. Is STRIPS a robot architecture?}\\
No. It is a symbolic planning method. An architecture must also organize perception, execution and the information exchanged among components.\par\end{samepage}

\begin{samepage}\textbf{11. What separates hierarchical and reactive sensing?}\\
Hierarchical sensing is fused into a shared world model for planning. Reactive sensing is distributed according to the needs of individual perception--action behaviors.\par\end{samepage}

\begin{samepage}\textbf{12. What does an affordance add to object recognition?}\\
It describes an opportunity for action relative to the agent. Detecting a passage wide enough for a robot can be more useful for navigation than naming every object around it.\par\end{samepage}

\begin{samepage}\textbf{13. Why combine deliberative and reactive layers?}\\
Deliberation supports goals, memory and future consequences; reactive behaviors support rapid responses to current conditions. An interface makes their decisions and status useful to each other.\par\end{samepage}

\begin{samepage}\textbf{14. Is collaborative robotics the same as a hybrid architecture?}\\
No. Collaboration concerns human--robot work. A hybrid architecture concerns the robot's internal organization of planning and behavior execution.\par\end{samepage}

\begin{samepage}\textbf{15. What are the three levels of a computational theory?}\\
Define the function, decompose it into inputs/outputs and transformations, then choose an implementation. Biological and robotic agents can share the first two without sharing the third.\par\end{samepage}

\begin{samepage}\textbf{16. What is an IRM, and what can be a releaser?}\\
An Innate Releasing Mechanism specifies activation of a behavior. A releaser can be an external stimulus, an internal state, or a combination, such as low battery and a visible station.\par\end{samepage}

\begin{samepage}\textbf{17. How do perceptual and motor schemas compose a behavior?}\\
The perceptual schema extracts the relevant information; the motor schema maps it into action. Instantiation supplies parameters for a concrete situation. The releaser determines activation.\par\end{samepage}

\begin{samepage}\textbf{18. How does STRIPS handle a failed precondition?}\\
In Murphy's means--ends presentation, the failed precondition becomes a subgoal. The planner recursively finds operators to achieve it, updates a modeled state, and returns to the original goal before physical execution.\par\end{samepage}

\begin{samepage}\textbf{19. Name the four criteria for evaluating an architecture.}\\
Support for modularity, niche targetability, ease of portability to other domains, and robustness. Explain each with respect to the intended application and its failure cases.\par\end{samepage}

\begin{samepage}\textbf{20. Does NHC repeat all planning after each observation?}\\
No. The Pilot monitors the current path segment; a reached waypoint advances navigation, while an obstruction can return control to the Navigator. The entire mission need not be recomputed.\par\end{samepage}

\begin{samepage}\textbf{21. Is active perception the same as an active sensor?}\\
No. Active perception changes an action to obtain better information. Whether the sensor hardware is active or passive is a separate classification (later sensing lessons).\par\end{samepage}

\begin{samepage}\textbf{22. Define SENSE, PLAN and ACT by their inputs and outputs.}\\
SENSE: sensor data $\rightarrow$ sensed information. PLAN: sensed and/or cognitive information $\rightarrow$ directives. ACT: sensed information or directives $\rightarrow$ actuator commands.\par\end{samepage}

\begin{samepage}\textbf{23. Describe the three paradigms by primitives and sensing organization.}\\
Hierarchical: S, P, A with sensing fused into a global world model. Reactive: concurrent S--A behaviors with local, behavior-specific sensing. Hybrid: P, S--A; sensing routed to behaviors and also to the planner's task-oriented world model.\par\end{samepage}

\begin{samepage}\textbf{24. Classical AI paradigm versus intelligent robotics paradigm?}\\
Classical AI reasons over an abstract, fully specified problem (chess) in a one-way perceive--reason--act chain. Intelligent robotics closes the loop through the real, partially known environment; autonomous behavior emerges from continuous perception--action interaction.\par\end{samepage}

\begin{samepage}\textbf{25. Why does robotics say ``self-governing'' without meaning moral independence?}\\
The term comes from the mechanical tradition (Watt's steam-engine governor): a machine regulating itself toward a goal. The dictionary sense of moral independence feeds the myth of evil robots; bounded rationality further limits any agent.\par\end{samepage}

\begin{samepage}\textbf{26. Give the four automation/autonomy keyword pairs of S04.}\\
Plans: execution vs generation. Actions: deterministic vs non-deterministic. Models: closed world vs open world. Knowledge representation: signals vs symbols. Automation is about tools, autonomy about agents.\par\end{samepage}

\begin{samepage}\textbf{27. Give two advantages and two disadvantages of the Hierarchical Paradigm.}\\
Advantages: a clear ordering of sensing, planning and acting; explicit goal reasoning (and a path from semi-autonomous to autonomous control). Disadvantages: the planning bottleneck on a global world model (frame problem, CWA); sensing and acting always disconnected, so no stimulus--response; uncertainty poorly handled.\par\end{samepage}

\begin{samepage}\textbf{28. Evaluate NHC/RCS with the four criteria.}\\
Modularity: intuitive but only functional decomposition. Niche targetability: reasonable (navigation; RCS on mining, submarines, cars). Portability: unclear, specified too broadly. Robustness: limited, via RCS plan simulation, at the price of delay.\par\end{samepage}

\subsection{One scenario across the whole course}
\textbf{Task:} a robot brings a component to a worker, but someone leaves a trolley in the corridor.

\textbf{Automation perspective.} If the design assumes a fixed clear route, replaying the prepared sequence may fail. A low-level feedback controller can track the requested motion without understanding that the route itself is unsuitable.

\textbf{Open-world perspective.} The prior map does not prove present clearance. The robot must detect the trolley and update its understanding of the route.

\textbf{Reactive response.} A local behavior slows or stops the robot in response to the obstruction. This response alone may leave the delivery unfinished.

\textbf{Deliberative response.} A planner evaluates an alternative route using the updated model and the delivery goal.

\textbf{Hybrid response.} The local behavior handles the immediate obstruction while the planning layer resolves the route. The interface connects the observed failure to the task plan.

\textbf{Shared-intelligence response.} If the route cannot be resolved, the robot can report the problem to a human. The human may change the goal or give contextual information while the robot performs the local execution.

\subsection{Common mistakes to avoid}
\begin{itemize}
\item Using \emph{closed world} and \emph{closed-loop control} interchangeably.
\item Treating an unknown fact as false in an open-world task.
\item Equating AI with machine learning, or intelligence with humanoid appearance.
\item Assuming autonomy means the complete absence of a human supervisor.
\item Assuming reactive behaviors do not need coordination or computation.
\item Assuming a world model is only a geometric map.
\item Using \emph{collaborative}, \emph{shared control} and \emph{hybrid architecture} as synonyms.
\end{itemize}
